% ============================================================
%  Modul 11 — Penjadwalan Proses Round-Robin
% ============================================================

\chapter{Penjadwalan Proses Round-Robin}
\label{chap:modul-11}

% --- 1. Tujuan dan Output Praktikum ---
\section{Tujuan dan Output Praktikum}

\subsection{Tujuan Praktikum}
Setelah menyelesaikan praktikum ini, mahasiswa diharapkan mampu:
\begin{enumerate}
	\item Memahami konsep \textit{CPU Scheduling} dan peran komponen \textit{scheduler} dalam mengelola pembagian giliran waktu CPU antar \textit{thread} pada tingkat kernel.
	\item Menganalisis perilaku algoritma penjadwalan \textit{Round-Robin} (FIFO) yang telah diimplementasikan pada Modul~8, khususnya pada fungsi \texttt{next\_thread\_to\_run()} di \texttt{src/thread.c}.
	\item Memahami konsep \textit{time slice}, \textit{cooperative scheduling}, dan implikasinya terhadap keadilan (\textit{fairness}) pembagian CPU.
	\item Menghubungkan perilaku \textit{scheduler} Round-Robin dengan keluaran (\textit{output}) sistem di QEMU dan membuktikan bahwa algoritma ini mengabaikan nilai prioritas \textit{thread}.
\end{enumerate}

\subsection{Output Praktikum}
Pada akhir praktikum ini, mahasiswa diharapkan menghasilkan:
\begin{itemize}
	\item Analisis kritis terhadap algoritma \textit{Round-Robin} yang berjalan di fungsi \texttt{next\_thread\_to\_run()}.
	\item Tangkapan layar (\textit{screenshot}) keluaran QEMU yang menunjukkan tiga \textit{thread} dengan prioritas berbeda (namun tetap berjalan bergantian secara merata) sebagai bukti bahwa \textit{Round-Robin} mengabaikan nilai prioritas.
	\item Laporan yang menjabarkan kelebihan dan keterbatasan \textit{Round-Robin}, khususnya ketidakmampuannya mendahulukan \textit{thread} berprioritas tinggi, sebagai landasan untuk modifikasi ke \textit{Priority Scheduling} pada Modul~12.
\end{itemize}

% --- 2. Dasar Teori ---
\newpage
\section{Dasar Teori}

\subsection{CPU Scheduling}
\textit{CPU Scheduling} adalah mekanisme fundamental yang digunakan sistem operasi untuk memutuskan \textit{thread} mana dari antrian \textit{ready} yang akan mendapat giliran eksekusi di CPU pada suatu waktu tertentu \cite{Silberschatz2018}. Komponen yang menjalankan keputusan ini disebut \textit{scheduler}. Tanpa \textit{scheduler}, CPU hanya mampu menjalankan satu alur eksekusi tanpa kemampuan beralih ke tugas lain.

Analoginya seperti sebuah lapangan sepak bola yang hanya memiliki satu bola (CPU). Meski ada 11 pemain (thread) di lapangan, pada satu waktu tertentu hanya satu pemain yang bisa menguasai bola. Pelatih (\textit{scheduler}) bertugas menentukan siapa yang harus menerima bola selanjutnya berdasarkan strategi tertentu.

Tujuan utama \textit{CPU scheduling} meliputi:
\begin{itemize}
	\item \textbf{Utilisasi CPU}: Memastikan CPU tidak pernah menganggur (\textit{idle}) selama masih ada \textit{thread} yang siap dijalankan. CPU yang menganggur sama dengan pemborosan sumber daya.
	\item \textbf{Fairness}: Memberikan kesempatan yang adil bagi setiap \textit{thread} untuk menggunakan CPU, sehingga tidak ada \textit{thread} yang terabaikan sepenuhnya.
	\item \textbf{Throughput}: Memaksimalkan jumlah \textit{thread} yang berhasil menyelesaikan pekerjaannya per satuan waktu.
	\item \textbf{Response Time}: Meminimalkan waktu tunggu antara permintaan \textit{thread} untuk dijalankan dan saat ia benar-benar mendapat giliran CPU.
\end{itemize}

\subsection{Algoritma Round-Robin (FIFO)}
\textit{Round-Robin} (RR) adalah algoritma penjadwalan yang paling sederhana dan paling intuitif. Nama ``Round-Robin'' berasal dari istilah yang digunakan dalam turnamen olahraga, di mana setiap peserta mendapat giliran bermain secara berputar.

Prinsipnya serupa dengan antrean di loket bank: siapa yang datang lebih dulu, ia yang dilayani lebih dulu (\textit{First-In, First-Out} / FIFO). Setelah mendapat giliran, \textit{thread} berjalan hingga salah satu dari dua hal terjadi:
\begin{enumerate}
	\item \textit{Thread} secara sukarela melepas CPU (melalui \texttt{thread\_yield()}).
	\item \textit{Thread} dihentikan paksa oleh \textit{timer interrupt} setelah \textit{time slice} habis (pada sistem \textit{preemptive}).
\end{enumerate}

Setelah melepas CPU, \textit{thread} tersebut masuk kembali ke belakang antrian, menunggu giliran berikutnya. Proses ini berulang terus sehingga setiap \textit{thread} mendapat giliran secara merata.

\subsubsection{Implementasi Ready Queue sebagai Circular Array}
Pada \textit{starter code} OS, antrian \textit{ready} diimplementasikan menggunakan \textit{circular array}:

\begin{lstlisting}[language=C, caption={Variabel antrian thread pada thread.c}]
static struct thread *ready_queue[MAX_THREADS * 2]; /* 16 slot */
static uint32_t ready_head  = 0;  /* Indeks depan antrian */
static uint32_t ready_tail  = 0;  /* Indeks belakang antrian */
static uint32_t ready_count = 0;  /* Jumlah thread di antrian */
\end{lstlisting}

Mekanisme \textit{circular array} bekerja sebagai berikut:
\begin{itemize}
	\item \textbf{Enqueue} (masukkan): Thread baru ditempatkan di posisi \texttt{ready\_tail}, lalu \texttt{ready\_tail} digeser maju satu posisi secara \textit{modular} (kembali ke 0 jika sudah mencapai ujung array).
	\item \textbf{Dequeue} (keluarkan): Thread diambil dari posisi \texttt{ready\_head}, lalu \texttt{ready\_head} digeser maju satu posisi.
\end{itemize}

Berikut ilustrasi keadaan antrian saat tiga thread (A, B, C) sedang terdaftar dan thread A baru saja dikeluarkan untuk dieksekusi:

\begin{lstlisting}[caption={Ilustrasi circular array ready\_queue}]
Indeks:  [0]   [1]   [2]   [3]   [4] ...
Isi:       -     B     C     -     -
                 ^           ^
          ready_head   ready_tail
          (dequeue)    (enqueue)
\end{lstlisting}

\subsection{Cooperative Scheduling vs Preemptive Scheduling}
Terdapat dua model utama dalam penjadwalan CPU:

\begin{enumerate}
	\item \textbf{Cooperative Scheduling}: \textit{Thread} secara sadar dan sukarela melepas CPU dengan memanggil fungsi seperti \texttt{thread\_yield()}. Jika sebuah \textit{thread} ``egois'' dan tidak pernah memanggil \texttt{yield()}, ia akan menguasai CPU selamanya, menyebabkan seluruh \textit{thread} lain membeku. Model inilah yang digunakan oleh IF-OS saat ini.

	\item \textbf{Preemptive Scheduling}: Sistem operasi menggunakan \textit{timer interrupt} dari perangkat keras untuk secara paksa merebut CPU dari \textit{thread} yang sedang berjalan setelah \textit{time slice} tertentu (misalnya 10 milidetik). Model ini digunakan oleh semua OS modern (Windows, Linux, macOS).
\end{enumerate}

\subsection{Field Priority pada Struct Thread}
Pada \texttt{src/thread.h}, setiap \texttt{struct thread} memiliki field bernama \texttt{priority}:

\begin{lstlisting}[language=C, caption={Definisi konstanta prioritas pada thread.h}]
#define PRI_MIN     0    /* Prioritas terendah */
#define PRI_DEFAULT 31   /* Prioritas default */
#define PRI_MAX     63   /* Prioritas tertinggi */
\end{lstlisting}

Field ini sudah diinisialisasi oleh \texttt{init\_thread()} setiap kali thread dibuat. Namun, fungsi \texttt{next\_thread\_to\_run()} pada implementasi \textit{Round-Robin} \textbf{sama sekali tidak membaca field ini}. Akibatnya, thread berprioritas 10 dan thread berprioritas 50 diperlakukan setara, keduanya mendapat giliran berdasarkan urutan kedatangan saja. Keterbatasan inilah yang akan diatasi pada Modul~12.

\subsection{Keterbatasan Round-Robin}
Meskipun \textit{Round-Robin} sangat adil (\textit{fair}), algoritma ini memiliki keterbatasan krusial:
\begin{itemize}
	\item \textbf{Tidak ada diferensiasi kepentingan:} Sebuah \textit{thread} yang menangani tugas kritis (misalnya penanganan \textit{interrupt} atau pemrosesan I/O) mendapat giliran yang sama dengan \textit{thread} yang hanya menjalankan tugas ringan.
	\item \textbf{Response time tidak optimal:} Thread penting yang baru masuk antrian harus menunggu semua thread di depannya selesai mendapat giliran, meskipun tugasnya lebih mendesak.
	\item \textbf{Tidak cocok untuk sistem real-time:} Sistem yang membutuhkan jaminan waktu respon (seperti kendali motor atau perangkat medis) tidak dapat mengandalkan Round-Robin.
\end{itemize}

Tabel \ref{tab:m11-dasar-teori} merangkum konsep-konsep penting pada modul ini.

\begin{table}[h]
	\caption{Tabel Ringkasan Dasar Teori Modul 11}
	\label{tab:m11-dasar-teori}
	\centering
	\begin{tabulary}{\textwidth}{|L|L|}
		\hline
		\textbf{Konsep} & \textbf{Penjelasan Singkat} \\ \hline
		CPU Scheduling
		& Mekanisme OS untuk memutuskan \textit{thread} mana yang mendapat giliran CPU berikutnya. \\ \hline
		Round-Robin (FIFO)
		& Algoritma penjadwalan berdasarkan urutan kedatangan; setiap \textit{thread} mendapat giliran secara merata tanpa memperhatikan prioritas. \\ \hline
		Time Slice / Quantum
		& Jatah durasi waktu CPU per giliran sebelum \textit{thread} harus menyerahkan kendali. \\ \hline
		Cooperative Scheduling
		& Model penjadwalan di mana \textit{thread} secara sukarela melepas CPU melalui pemanggilan \texttt{yield()}. \\ \hline
		Preemptive Scheduling
		& Model penjadwalan di mana OS memaksa perpindahan \textit{thread} melalui \textit{timer interrupt}. \\ \hline
		Ready Queue (Circular Array)
		& Antrian penampung \textit{thread} berstatus \texttt{THREAD\_READY}; dikelola dengan indeks \texttt{ready\_head} dan \texttt{ready\_tail}. \\ \hline
	\end{tabulary}
\end{table}

% --- 3. Alat dan Bahan ---
\newpage
\section{Alat dan Bahan}
\label{sec:m11-alat-bahan}

\subsection{Perangkat Keras}
\begin{itemize}
	\item PC atau Laptop dengan arsitektur prosesor basis \textit{host} x86/x64 yang mendukung kelanjutan kompilasi emulator.
	\item Memori minimal 4GB RAM ruang eksekusi.
\end{itemize}

\subsection{Perangkat Lunak (Bahan Praktikum)}
\begin{itemize}
	\item OS \textit{Host} utama (Windows dengan WSL/GitBash, Linux, atau macOS).
	\item Text Editor (misalnya Visual Studio Code) untuk menyunting berkas kode bahasa C dan \textit{Assembly}.
	\item Instalasi Docker Engine (\texttt{docker.io}) untuk menyalakan lingkungan isolasi kompilasi.
	\item Emulator Hardware x86: \texttt{qemu-system-x86} (atau instalasi QEMU setara) yang akan mengeksekusi binari hasil OS.
	\item Direktori \textit{Starter Code} OS 32-bit dengan implementasi \texttt{thread\_create()}, \texttt{thread\_yield()}, dan \texttt{next\_thread\_to\_run()} yang sudah berfungsi dari Modul~8 (berisi direktori \texttt{src} dengan kerangka lengkap \texttt{kernel.c}, \texttt{thread.h}, \texttt{thread.c}, \texttt{switch.asm}, dsb).
\end{itemize}

% --- 4. Langkah-Langkah Praktikum ---
\newpage
\section{Langkah-Langkah Praktikum}

\subsection{Persiapan Lingkungan}
\label{sec:m11-persiapan}

Sebelum memulai praktikum, pastikan \textit{toolchain} dan \textit{source code} OS sudah tersedia.
Karena praktikum ini merupakan lanjutan dari implementasi \textit{multi-threading} (Modul~8) dan sinkronisasi (Modul~9--10), implementasi \texttt{thread\_create()}, \texttt{thread\_yield()}, dan \texttt{next\_thread\_to\_run()} di \texttt{thread.c} \textbf{harus sudah selesai dan berjalan dengan benar}. Jika Anda menggunakan \textit{starter code} mentah (tanpa implementasi Modul~8), simulasi akan terhenti seketika dan menghasilkan pesan \textit{error}: \texttt{kernel\_main() kembali dari thread\_block()!}.

\begin{enumerate}
	\item Pastikan Anda sudah memiliki direktori \textit{starter code} OS dengan file \texttt{src/thread.c} yang sudah diimplementasikan pada Modul~8 (atau modul sebelumnya). Jika Anda sudah memiliki direktori kerja dari Modul~9 atau Modul~10 yang berfungsi, Anda dapat langsung menggunakannya.

	      \textit{Catatan: Jika Anda belum memiliki direktori kerja, lakukan langkah berikut untuk membuat yang baru:}
	      \begin{lstlisting}[language=bash, caption={Membuat direktori kerja baru (opsional, jika belum punya)}]
mkdir praktikum-minggu6
cd praktikum-minggu6
git clone https://github.com/sinavarasina/starter-code-os32.git
cd starter-code-os32
    \end{lstlisting}

	\item \textbf{Langkah Kritis (hanya jika clone baru):} Jika Anda baru saja melakukan \textit{clone} ulang repositori pada langkah sebelumnya, file \texttt{src/thread.c} masih berupa \textit{stub} kosong. Anda \textbf{harus} menimpanya dengan implementasi \texttt{thread.c} yang sudah berfungsi dari Modul 8.

	      \textbf{Cara Salin manual:} Salin file \texttt{thread.c} yang sudah Anda kerjakan pada Modul 8 ke dalam folder \texttt{src/}.
	      \begin{lstlisting}[language=bash, caption={Menyalin thread.c dari direktori sebelumnya}]
cp /path/ke/modul8/starter-code-os32/src/thread.c src/thread.c
    \end{lstlisting}
	      \textit{Catatan: Jika Anda tidak menyimpan file \texttt{thread.c} dari pertemuan sebelumnya, silakan hubungi asisten praktikum.}

	      \textbf{Penting:} File \texttt{thread.c} yang disalin harus mengandung baris \texttt{\#include "thread.h"} di awal file. Jika baris tersebut tidak ada, tambahkan tiga baris berikut di \textbf{paling atas} file \texttt{src/thread.c}:
	      \begin{lstlisting}[language=C, caption={Header include yang wajib ada di awal thread.c}]
#include "thread.h"
#include "types.h"
#include "vga.h"
    \end{lstlisting}

	\item Verifikasi bahwa seluruh berkas sumber (\textit{source code}) sudah lengkap.
	      \begin{lstlisting}[language=bash, caption={Verifikasi direktori source code}]
ls src/
# Harus menampilkan: boot.asm  kernel.c  scheduler.h  synch.c  synch.h
#                    switch.asm  thread.c  thread.h
#                    types.h  vga.c  vga.h
    \end{lstlisting}

	      \textit{Jika ada file yang hilang, ulangi proses clone dari langkah 1.}
\end{enumerate}

% -----------------------------------------------------------------
\subsection{Sesi 1: Menjalankan Demo Round-Robin dengan Thread Berprioritas Berbeda}
\label{sec:m11-sesi-1}

Pada sesi ini, praktikan akan mengompilasi dan menjalankan OS dengan \textit{flag} \texttt{PRIORITY\_SCHEDULING}.
\textbf{Tidak perlu mengubah kode apa pun}, cukup amati keluaran (\textit{output}) yang muncul di layar QEMU.

\textbf{Penting:} Flag \texttt{PRIORITY\_SCHEDULING} pada Makefile \textbf{hanya mengubah skenario demo} di \texttt{kernel.c}, yaitu membuat tiga \textit{thread} dengan prioritas berbeda (\texttt{Low}, \texttt{Med}, \texttt{High}). Flag ini \textbf{tidak} mengubah algoritma \textit{scheduler}. Algoritma \textit{scheduler} tetap ditentukan oleh implementasi fungsi \texttt{next\_thread\_to\_run()} di \texttt{thread.c}, yang saat ini masih berupa \textit{Round-Robin} (FIFO).

Perhatikan bahwa pada file \texttt{kernel.c}, blok \texttt{\#elif defined(PRIORITY\_SCHEDULING)} membuat tiga \textit{thread} dengan prioritas \textbf{berbeda}:

\begin{lstlisting}[language=C, caption={Pembuatan thread dengan prioritas berbeda di kernel.c}]
thread_create("Low",  10,          thread_func_generic, &args_low);
thread_create("Med",  PRI_DEFAULT, thread_func_generic, &args_med);
thread_create("High", 50,          thread_func_generic, &args_high);
\end{lstlisting}

Konfigurasi argumen masing-masing thread:
\begin{itemize}
	\item \texttt{Low}: prioritas 10, warna abu-abu (\texttt{VGA\_COLOR\_DARK\_GRAY}), berjalan tanpa batas (\texttt{max\_ticks = 0}).
	\item \texttt{Med}: prioritas 31 (\texttt{PRI\_DEFAULT}), warna cyan terang (\texttt{VGA\_COLOR\_BRIGHT\_CYAN}), berjalan tanpa batas.
	\item \texttt{High}: prioritas 50, warna putih terang (\texttt{VGA\_COLOR\_BRIGHT\_WHITE}), berhenti setelah 5 tick (\texttt{max\_ticks = 5}).
\end{itemize}

\begin{enumerate}
	\item \textit{Aktivasi lingkungan tersetup \& terisolasi}.
	      \begin{lstlisting}[language=bash, caption={Menjalankan Docker menggunakan script}]
./run-docker.sh
    \end{lstlisting}

	      \textit{Catatan: Jika muncul error \texttt{permission denied}, jalankan \texttt{chmod +x run-docker.sh} terlebih dahulu. Jika Docker belum berjalan, pastikan Docker Desktop atau Docker Engine sudah aktif.}

	\item \textit{Compile} OS dengan \textit{flag} \texttt{PRIORITY\_SCHEDULING} diaktifkan.
	      \begin{lstlisting}[language=bash, caption={Compile dengan flag PRIORITY\_SCHEDULING}]
make clean
make DEMO=-DPRIORITY_SCHEDULING
    \end{lstlisting}

	      \textbf{PERHATIAN (Kesalahan Umum):}
	      \begin{itemize}
	      	\item Perintah \texttt{make clean} \textbf{wajib} dijalankan terlebih dahulu untuk menghapus hasil kompilasi sebelumnya. Jika langkah ini dilewati, kode lama yang sudah di-compile mungkin masih digunakan dan flag \texttt{PRIORITY\_SCHEDULING} tidak akan terdeteksi.
	      	\item Jika muncul error \texttt{thread.c: undefined reference to `allocate\_tid'}, berarti file \texttt{thread.c} yang digunakan tidak sesuai atau belum lengkap. Pastikan file mengandung baris \texttt{\#include "thread.h"} di awal, dan ulangi langkah Persiapan.
	      \end{itemize}

	\item Jalankan kernel di QEMU dan amati keluaran (\textit{output}) di layar.
	      \begin{lstlisting}[language=bash, caption={Menjalankan kernel di QEMU}]
make run
    \end{lstlisting}

	      \textit{Jika muncul error \texttt{qemu-system-i386: command not found}, QEMU belum terinstal pada mesin host Anda. Keluar dari Docker (ketik \texttt{exit}), instal QEMU di host, lalu jalankan \texttt{make run} dari terminal host.}

	\item Amati keluaran yang muncul selama kurang lebih 10--15 detik, ambil tangkapan layar (\textit{screenshot}), lalu tutup QEMU (tekan \texttt{q} atau \texttt{Ctrl+C}).

	      Karena \texttt{next\_thread\_to\_run()} masih menggunakan FIFO (\textit{Round-Robin}), ketiga \textit{thread} bergantian \textbf{tanpa memperhatikan prioritas}:

	      \begin{lstlisting}[caption={Keluaran QEMU: Round-Robin mengabaikan prioritas}]
===================================
  IF-OS  |  Praktikum Sistem Operasi
===================================

Inisialisasi sistem thread...
Demo Minggu 6: Priority Scheduling
(Jika masih Round-Robin, next_thread_to_run() belum dimodifikasi)
(High=50 jalan 5 tick lalu exit, Med=31 mengambil alih, Low=10 terakhir)

[Low] tick 0
[Med] tick 0
[High] tick 0
[Low] tick 1
[Med] tick 1
[High] tick 1
...
[High] selesai (5 tick)
[Low] tick 5
[Med] tick 5
[Low] tick 6
[Med] tick 6
...
    \end{lstlisting}

	      \textit{Catatan: Pesan header pada output QEMU mendeskripsikan perilaku yang baru akan terjadi setelah Modul 12. Pada modul ini, ketiga thread masih berjalan bergantian secara merata.}

	      Perhatikan bahwa \texttt{Low} (prioritas 10) mendapat giliran yang \textbf{sama} dengan \texttt{High} (prioritas 50). Ketiganya bergilir secara merata: Low $\to$ Med $\to$ High $\to$ Low $\to$ \ldots{}. Dalam skenario sistem nyata, ini adalah perilaku yang tidak diinginkan karena tugas berprioritas tinggi seharusnya didahulukan.

	\item Bandingkan keluaran di atas dengan tabel ekspektasi berikut:

	      \begin{table}[htbp]
		      \centering
		      \caption{Ekspektasi urutan eksekusi pada Round-Robin}
		      \begin{tabulary}{\textwidth}{|c|l|c|L|}
			      \hline
			      \textbf{Giliran} & \textbf{Thread} & \textbf{Prioritas} & \textbf{Keterangan} \\ \hline
			      1 & Low  & 10 & Masuk antrian pertama kali \\ \hline
			      2 & Med  & 31 & Masuk antrian kedua        \\ \hline
			      3 & High & 50 & Masuk antrian ketiga        \\ \hline
			      4 & Low  & 10 & Kembali ke depan antrian (FIFO berputar) \\ \hline
			      5 & Med  & 31 & \ldots{} dan seterusnya   \\ \hline
		      \end{tabulary}
	      \end{table}

	      Urutan di atas membuktikan bahwa \textit{scheduler} \textit{Round-Robin} sepenuhnya mengabaikan kolom ``Prioritas''. Thread Low (10) selalu dieksekusi sebelum High (50) karena Low masuk antrian lebih dulu.
\end{enumerate}

\begin{figure}[h]
	\centering
	% TODO: ganti dengan screenshot hasil QEMU Round-Robin
	\includegraphics[width=0.85\textwidth]{Figure/11/1.png}
	\caption{Keluaran QEMU menunjukkan ketiga thread berjalan bergantian secara merata meskipun prioritasnya berbeda, yang membuktikan bahwa Round-Robin mengabaikan nilai prioritas. Keterbatasan ini akan diatasi pada Modul~12.}
	\label{fig:m11-sesi1}
\end{figure}
